\section{Conclusion}

Approximate Sign--Ratio Recovery extracts a shared assignment by
clustering profiles of aligned and normalized source weights. The
assignment fixes the sharing groups on each new task, which learns
its own numerical values.

For aligned source weights with fixed active connections, we characterize
exact recovery through scale-invariant magnitudes and sign profiles.
A condition on the graph of co-occurring codes gives uniqueness even
when no pair of codes repeats across neurons.
For partial observations, a completion criterion and a functional
ReLU example show how several individually ambiguous tasks can
determine the groups jointly. Recovery from erroneous pairwise
relations additionally requires sufficient redundancy between
the observed parts.

With known support, aSRR recovers the signed groups in all six
cross-correlation layouts. On four of five real datasets, the learned
assignments improve accuracy over shuffles at the same coefficient
budget. The continuous bases of MSR-FC and Yeh (soft) attain higher
accuracy in several settings, including FC100, WISDM and transformed
Aug-Omniglot. The auxiliary sparse ReLU experiment illustrates the
importance of the candidate support. For that teacher, the aSRR bank
retains all edges, and its selected assignment has higher query error
than Dense.
